% ============================================================
% SECTION 6: ISLAMIC GOLDEN AGE MATHEMATICS (~800 - 1400 CE)
% Slides 37-50 (14 slides, 11 mathematicians)
% This is intentionally the LARGEST section.
% ============================================================

% ---- Section Divider ----
{
\setbeamercolor{background canvas}{bg=islamicGreen}
\begin{frame}[plain]
  \centering
  \vfill
  \textcolor{white}{\Huge\bfseries The House of Wisdom}\\[0.7em]
  \textcolor{white!85}{\Large When Baghdad Was the Center of the World}\\[1.5em]
  \textcolor{white!65}{\large 800--1400 CE}
  \vfill
  % Timeline marker
  \visualplaceholder[5cm]{Diagram}
  \vfill
\end{frame}
}

\section{The Islamic Golden Age}

% ---- Slide 38: The Translation Movement ----
\begin{frame}{The Translation Movement}
  \begin{columns}[T]
    \column{0.55\textwidth}
    \begin{itemize}
      \item Under the Abbasid Caliphate, a massive translation project preserved Greek, Indian, and Persian texts in Arabic
      \item The \textbf{House of Wisdom} (\textit{Bayt al-Hikma}) in Baghdad was the epicenter
      \item Without this effort, much of Greek mathematics might have been permanently lost
    \end{itemize}
    \vspace{0.5em}
    \begin{insightbox}
      When Europeans later ``rediscovered'' Greek mathematics, they were often translating \emph{Arabic} translations of Greek originals that no longer existed in Greek.
    \end{insightbox}

    \column{0.42\textwidth}
    % Knowledge transmission flow chart
    \visualplaceholder[5cm]{Greek Texts  (some lost originals)}
  \end{columns}

  \note{
    \textbf{Talking point:} ``When Europeans later `rediscovered' Greek mathematics,
    they were often translating Arabic translations of Greek originals that no longer
    existed in Greek.''

    \verifytag{The House of Wisdom's exact nature is debated. Some historians (Gutas)
    describe it primarily as a library/translation center; others see it as more of an
    academy. Present cautiously --- say ``center of learning'' rather than ``university.''}
  }
\end{frame}

% ---- Slide 39: Al-Khwarizmi ----
\begin{frame}{Muhammad ibn Musa al-Khwarizmi --- The Father of Algebra}
  \begin{columns}[T]
    \column{0.58\textwidth}
    \begin{personbox}[Al-Khwarizmi]{islamicGreen}
      \textbf{Full name:} Abu Abdallah Muhammad ibn Musa al-Khwarizmi\\
      % Arabic: ابو عبد الله محمد بن موسى الخوارزمي
      {\small (Arabic name --- see manuscript image)}\\[3pt]
      \textbf{Dates:} c.\ 780 -- c.\ 850 CE\\
      \textbf{Origin:} Khwarezm (modern Khiva, Uzbekistan);\\
      worked in Baghdad at the House of Wisdom
    \end{personbox}
    \vspace{0.4em}
    \begin{itemize}
      \item Wrote \textit{Kitab al-Jabr wa-l-Muqabala} (c.\ 820)
      \item The word \textbf{``algebra''} comes from \textit{al-jabr} (``completion'')
      \item The word \textbf{``algorithm''} comes from the Latin form of his name
      \item Introduced the Hindu-Arabic decimal system to the Islamic world
    \end{itemize}

    \column{0.39\textwidth}
    % Etymology visual
    \visualplaceholder[5cm]{Arabic script: al-jabr -- requires Arabic font image}

    \vspace{0.5em}
    \visualplaceholder[3.5cm]{Page from the \textit{al-Jabr} manuscript\\(Bodleian Library MS.\@ Huntington 214)}
  \end{columns}

  \note{
    \textbf{Talking point:} ``Two words you use constantly --- `algebra' and `algorithm'
    --- both come from this one mathematician.''

    \verifytag{Al-Khwarizmi's birth year is uncertain: some sources say c.\ 780, others c.\ 790.}
  }
\end{frame}

% ---- Slide 40: Al-Khwarizmi's Algebra in Detail ----
\begin{frame}{Al-Khwarizmi's Algebra: Completing the Square}
  \begin{columns}[T]
    \column{0.48\textwidth}
    {\small\textbf{His words (paraphrased):}}\\[2pt]
    {\footnotesize\itshape ``A square and ten roots equal thirty-nine dirhams''}\\[4pt]
    {\small Modern form: $x^2 + 10x = 39$}

    \vspace{0.6em}
    {\small\textbf{Geometric method:}}

    % Completing the square -- geometric TikZ
    \visualplaceholder[5cm]{Completing the square -- geometric diagram}

    \column{0.48\textwidth}
    {\small\textbf{Step-by-step:}}
    \begin{enumerate}\setlength\itemsep{2pt}
      \item[\textcolor{islamicGreen}{\textbf{1.}}] {\footnotesize Start with $x^2 + 10x = 39$}
      \item[\textcolor{islamicGreen}{\textbf{2.}}] {\footnotesize Draw the $x^2$ square}
      \item[\textcolor{islamicGreen}{\textbf{3.}}] {\footnotesize Distribute $10x$ as two rectangles of $\frac{5}{2} \times x$}
      \item[\textcolor{islamicGreen}{\textbf{4.}}] {\footnotesize \textcolor{accentOrange}{\textbf{Complete}} the square: add $\left(\frac{5}{2}\right)^2 = \frac{25}{4}$}
      \item[\textcolor{islamicGreen}{\textbf{5.}}] {\footnotesize Now: $\left(x+\frac{5}{2}\right)^2 = 39 + \frac{25}{4} = \frac{181}{4}$}
      \item[\textcolor{islamicGreen}{\textbf{6.}}] {\footnotesize $x + \frac{5}{2} = \frac{\sqrt{181}}{2}$, so $x = \frac{\sqrt{181}-5}{2} = 3$}
    \end{enumerate}

    \vspace{0.5em}
    \begin{insightbox}
      {\footnotesize He solved the same equations you solve --- but geometrically, in words, with no symbols.}
    \end{insightbox}
  \end{columns}

  \note{
    \textbf{Talking point:} ``He solved the same equations you solve in class --- but
    instead of $x^2 + 10x = 39$, he wrote: `A square and ten roots equal thirty-nine
    dirhams.'\,''

    Note: $x^2 + 10x = 39$ gives $x = 3$ exactly. Check: $9+30=39$. Correct.

    Al-Khwarizmi classified six types of equations and gave geometric solutions for each.
    This geometric completing the square is the most famous.
  }
\end{frame}

% ---- Slide 41: Thabit ibn Qurra ----
\begin{frame}{Thabit ibn Qurra --- Translator and Innovator}
  \begin{columns}[T]
    \column{0.55\textwidth}
    \begin{personbox}[Thabit ibn Qurra]{islamicGreen}
      % Arabic: ثابت بن قرة
      {\small (Thabit ibn Qurra)}\\[3pt]
      \textbf{Dates:} 826 -- 901 CE\\
      \textbf{Origin:} Harran (modern Turkey); worked in Baghdad
    \end{personbox}
    \vspace{0.3em}
    \begin{itemize}
      \item Prolific translator: Euclid, Archimedes, Apollonius
      \item Extended Euclid's number theory
      \item Proved a generalization of the Pythagorean theorem to arbitrary triangles
      \item Discovered a rule for finding \textbf{amicable numbers}
    \end{itemize}

    \vspace{0.3em}
    \begin{insightbox}
      {\footnotesize \textbf{Amicable numbers:} 220 and 284. Each equals the sum of the other's proper divisors.\\[2pt]
      $220 \to 1{+}2{+}4{+}5{+}10{+}11{+}20{+}22{+}44{+}55{+}110 = 284$\\
      $284 \to 1{+}2{+}4{+}71{+}142 = 220$}
    \end{insightbox}

    \column{0.42\textwidth}
    % Generalized Pythagorean theorem diagram
    \visualplaceholder[5cm]{Arbitrary triangle}

    \vspace{0.5em}
    \visualplaceholder[3.5cm]{Manuscript page from Thabit's translation of Euclid}
  \end{columns}

  \note{
    \textbf{Key point:} Thabit was not just a translator --- he was a creative mathematician
    who extended the work he translated.

    The amicable numbers 220 and 284 were known to the Pythagoreans, but Thabit discovered
    a \emph{general rule} for finding such pairs.

    \verifytag{Do NOT present a specific formula for Thabit's amicable number rule ---
    describe qualitatively only, as the exact formulation needs verification from primary
    sources (Rashed's edition or Dickson's ``History of the Theory of Numbers'').}
  }
\end{frame}

% ---- Slide 42: Abu Kamil ----
\begin{frame}{Abu Kamil --- Extending the Algebra}
  \begin{columns}[T]
    \column{0.55\textwidth}
    \begin{personbox}[Abu Kamil Shuja ibn Aslam]{islamicGreen}
      % Arabic: ابو كامل شجاع بن اسلم
      {\small (Abu Kamil Shuja ibn Aslam)}\\[3pt]
      \textbf{Dates:} c.\ 850 -- c.\ 930 CE\\
      \textbf{Origin:} Egypt
    \end{personbox}
    \vspace{0.4em}
    \begin{itemize}
      \item Extended al-Khwarizmi's algebra \emph{significantly}
      \item Worked with \textbf{irrational numbers} as both solutions \emph{and} coefficients
        \begin{itemize}
          \item[\footnotesize$\bullet$] {\footnotesize Al-Khwarizmi had restricted to rationals}
        \end{itemize}
      \item Solved systems of equations
      \item Influenced al-Karaji and, through Latin translation, \textbf{Fibonacci}
    \end{itemize}

    \column{0.42\textwidth}
    \visualplaceholder[3.8cm]{A problem from Abu Kamil's \textit{Algebra} with its geometric solution}

    \vspace{0.5em}
    % Influence chain
    \visualplaceholder[5cm]{Al-Khwarizmi}
  \end{columns}

  \note{
    Abu Kamil is sometimes called ``the Egyptian calculator.'' He is the crucial link
    between al-Khwarizmi and both the later Islamic algebraic tradition (al-Karaji)
    and European mathematics (Fibonacci).

    \verifytag{Exact dates uncertain. Some sources give c.\ 850--c.\ 930, others c.\ 850--c.\ 935.}
  }
\end{frame}

% ---- Slide 43: Al-Karaji ----
\begin{frame}{Al-Karaji --- Algebra Freed from Geometry}
  \begin{columns}[T]
    \column{0.55\textwidth}
    \begin{personbox}[Abu Bakr al-Karaji]{islamicGreen}
      % Arabic: ابو بكر الكرجي
      {\small (sometimes written al-Karkhi)}\\[3pt]
      \textbf{Dates:} c.\ 953 -- c.\ 1029 CE\\
      \textbf{Origin:} Karaj (or Karkh, Baghdad); worked in Baghdad
    \end{personbox}
    \vspace{0.3em}
    \begin{itemize}
      \item Freed algebra from geometry --- treated algebraic operations \textbf{purely arithmetically}
      \item Proved the \textbf{binomial theorem}
      \item Proved the sum-of-cubes formula using a method akin to \textbf{mathematical induction}
      \item Extended algebra to powers beyond cubes
    \end{itemize}

    \column{0.42\textwidth}
    % Sum of cubes visual proof
    \visualplaceholder[5cm]{Visual representation: stacked squares}

    \vspace{0.6em}
    \begin{insightbox}
      {\footnotesize He arguably used the first form of \textbf{mathematical induction} --- a proof technique you may encounter at university.}
    \end{insightbox}
  \end{columns}

  \note{
    \textbf{Talking point:} ``He was arguably the first person to use something like
    mathematical induction --- a proof technique you may encounter in university.''

    Al-Karaji's key insight: algebraic manipulations do not \emph{need} geometric
    justification. This was a conceptual leap.

    \verifytag{Whether his method truly constitutes ``mathematical induction'' is debated.
    It is proto-inductive but not formalized as such. Say ``akin to'' or ``anticipating.''}
  }
\end{frame}

% ---- Slide 44: Omar Khayyam ----
\begin{frame}{Omar Khayyam --- Poet and Mathematician}
  \begin{columns}[T]
    \column{0.53\textwidth}
    \begin{personbox}[Omar Khayyam]{islamicGreen}
      {\small Ghiyath al-Din Abu l-Fath Umar ibn Ibrahim al-Khayyam\\
      % Arabic: غياث الدين ابو الفتح عمر بن ابراهيم الخيام
      }\\[3pt]
      \textbf{Dates:} 1048 -- 1131 CE\\
      \textbf{Origin:} Nishapur, Khorasan (modern Iran)
    \end{personbox}
    \vspace{0.3em}
    \begin{itemize}
      \item Classified and solved \textbf{all cubic equations} geometrically (intersecting conic sections)
      \item Contributed to the \textbf{Jalali calendar} reform
      \item Famous for the \textit{Rubaiyat} poetry
    \end{itemize}

    \column{0.44\textwidth}
    % Conic section construction for cubics (pure TikZ)
    \visualplaceholder[5cm]{Parabola diagram}
  \end{columns}

  \note{
    \textbf{Talking point:} ``He is one of the few people equally famous as a mathematician
    AND a poet. His calendar was accurate to within one day in 5,000 years.''

    The title \textit{al-Jabr wa l-Muqabala} was a standard genre name for algebra treatises,
    not unique to Khayyam or al-Khwarizmi.

    \verifytag{Jalali calendar accuracy: standard claim is 1 day in $\sim$5,000 years vs.\
    Gregorian's 1 day in $\sim$3,236 years. Verify these specific numbers before presenting.}
  }
\end{frame}

% ---- Slide 45: Ibn al-Haytham ----
\begin{frame}{Ibn al-Haytham --- The First Scientist?}
  \begin{columns}[T]
    \column{0.55\textwidth}
    \begin{personbox}[Ibn al-Haytham (Alhazen)]{islamicGreen}
      {\small Abu Ali al-Hasan ibn al-Hasan ibn al-Haytham\\
      % Arabic: ابو علي الحسن بن الحسن بن الهيثم
      known in Latin as Alhazen}\\[3pt]
      \textbf{Dates:} c.\ 965 -- c.\ 1040 CE\\
      \textbf{Origin:} Basra (modern Iraq); worked in Cairo
    \end{personbox}
    \vspace{0.3em}
    \begin{itemize}
      \item \textit{Kitab al-Manazir} (Book of Optics) --- revolutionized the understanding of vision and light
      \item Solved \textbf{Alhazen's problem}: finding the reflection point on a spherical mirror (4th-degree equation)
      \item Used an early form of \textbf{integral calculus} to compute the volume of a paraboloid
      \item Pioneered the \textbf{scientific method}
    \end{itemize}

    \column{0.42\textwidth}
    % Alhazen's problem diagram
    \visualplaceholder[5cm]{Spherical mirror}

    \vspace{0.4em}
    \visualplaceholder[3.5cm]{Diagram from the \textit{Book of Optics}\\(many historical copies survive)}
  \end{columns}

  \note{
    \textbf{Talking point:} ``Some historians call him the `first true scientist' because
    he insisted on experimental verification --- hypothesis, experiment, conclusion.''

    His Book of Optics influenced Roger Bacon, Kepler, and Descartes. The Latin translation
    (\textit{De Aspectibus}) was widely read in medieval Europe.
  }
\end{frame}

% ---- Slide 46: Al-Biruni ----
\begin{frame}{Al-Biruni --- The Polymath \hfill {\small\textcolor{islamicGreen!60}{[EXTENDED]}}}
  \begin{columns}[T]
    \column{0.50\textwidth}
    \begin{personbox}[Abu Rayhan al-Biruni]{islamicGreen}
      % Arabic: ابو ريحان محمد بن احمد البيروني
      {\small (Abu Rayhan al-Biruni)}\\[3pt]
      \textbf{Dates:} 973 -- c.\ 1048 CE\\
      \textbf{Origin:} Kath, Khwarezm (modern Uzbekistan)
    \end{personbox}
    \vspace{0.3em}
    \begin{itemize}
      \item Calculated the Earth's radius to within \textbf{$\sim$1\%} of the modern value
      \item Used trigonometry and a mountain
      \item Wrote \textit{Tarikh al-Hind} on Indian mathematics and astronomy
      \item Pioneered comparative anthropology and geodesy
    \end{itemize}

    \column{0.47\textwidth}
    % Earth radius measurement diagram
    \visualplaceholder[5cm]{Earth (partial arc)}
  \end{columns}

  \note{
    \textbf{Talking point:} ``He calculated the size of the Earth using a mountain and
    trigonometry --- and got it right to within about 15 kilometers.''

    Al-Biruni measured the angle of dip to the horizon from the top of a mountain
    of known height, then used trigonometry to compute the Earth's radius.

    \verifytag{Exact accuracy claim needs verification. ``Within 1\%'' and ``within 15~km''
    are commonly cited but should be checked against the modern value of $\sim$6,371~km.}
  }
\end{frame}

% ---- Slide 47: Al-Samawal and Sharaf al-Din al-Tusi ----
\begin{frame}{Al-Samawal \& Sharaf al-Din al-Tusi \hfill {\small\textcolor{islamicGreen!60}{[EXTENDED]}}}
  \begin{columns}[T]
    \column{0.48\textwidth}
    \begin{personbox}[Al-Samawal]{islamicGreen}
      % Arabic: السموأل بن يحيى المغربي
      {\small (al-Samawal ibn Yahya al-Maghribi)}\\[2pt]
      \textbf{Dates:} c.\ 1130 -- c.\ 1180 CE\\
      \textbf{Origin:} Baghdad
    \end{personbox}
    \vspace{0.2em}
    \begin{itemize}\setlength\itemsep{1pt}
      \item Advanced algebraic notation
      \item First statement of \textbf{polynomial long division}
    \end{itemize}

    \vspace{0.4em}
    % Polynomial division example
    \visualplaceholder[5cm]{Diagram}

    \column{0.48\textwidth}
    \begin{personbox}[Sharaf al-Din al-Tusi]{islamicGreen}
      % Arabic: شرف الدين الطوسي
      {\small (Sharaf al-Din al-Tusi)}\\[2pt]
      \textbf{Dates:} c.\ 1135 -- 1213 CE\\
      \textbf{Origin:} Tus (modern Iran)
    \end{personbox}
    \vspace{0.2em}
    \begin{itemize}\setlength\itemsep{1pt}
      \item Analyzed cubic equations by finding their \textbf{maxima and minima}
      \item Anticipating differential calculus!
      \item \textit{Not} the same person as Nasir al-Din al-Tusi
    \end{itemize}

    \vspace{0.2em}
    % Cubic curve with max/min (pure TikZ)
    \visualplaceholder[5cm]{Cubic equation analysis diagram}
  \end{columns}

  \note{
    These two figures show the continued vitality of Islamic mathematics well into the
    12th--13th centuries.

    Sharaf al-Din al-Tusi's analysis of maxima/minima of cubics is remarkable ---
    it is essentially the same idea as using the derivative to find critical points,
    centuries before calculus was formalized.

    \verifytag{Al-Samawal's dates uncertain. Also verify the polynomial division attribution.}

    \textbf{Important:} Sharaf al-Din al-Tusi ($\sim$1135--1213) is NOT the same person as
    Nasir al-Din al-Tusi (1201--1274). Make this distinction explicit.
  }
\end{frame}

% ---- Slide 48: Nasir al-Din al-Tusi ----
\begin{frame}{Nasir al-Din al-Tusi --- Trigonometry Becomes Its Own Field}
  \begin{columns}[T]
    \column{0.53\textwidth}
    \begin{personbox}[Nasir al-Din al-Tusi]{islamicGreen}
      {\small Muhammad ibn Muhammad ibn al-Hasan al-Tusi\\
      % Arabic: محمد بن محمد بن الحسن الطوسي
      known as Nasir al-Din al-Tusi}\\[3pt]
      \textbf{Dates:} 1201 -- 1274 CE\\
      \textbf{Origin:} Tus, Khorasan (modern Iran)
    \end{personbox}
    \vspace{0.3em}
    \begin{itemize}
      \item Wrote \textit{Treatise on the Quadrilateral} --- the \textbf{first work} to treat trigonometry as an independent discipline
      \item Stated and proved the \textbf{law of sines} for plane and spherical triangles
      \item Built the \textbf{Maragheh observatory}, one of the most advanced of its era
    \end{itemize}

    \column{0.44\textwidth}
    % Law of sines illustration
    \visualplaceholder[5cm]{Triangle}

    \vspace{0.3em}
    \visualplaceholder[3.5cm]{Reconstruction of the Maragheh observatory}
  \end{columns}

  \note{
    \textbf{Talking point:} ``Before al-Tusi, trigonometry was just a tool for astronomy.
    He made it its own branch of mathematics.''

    The Maragheh observatory was built after the Mongol conquest, under the patronage of
    Hulagu Khan. It shows how mathematical activity continued even through periods of
    political upheaval.
  }
\end{frame}

% ---- Slide 49: Jamshid al-Kashi ----
\begin{frame}{Jamshid al-Kashi --- Master of Computation}
  \begin{columns}[T]
    \column{0.55\textwidth}
    \begin{personbox}[Jamshid al-Kashi]{islamicGreen}
      {\small Ghiyath al-Din Jamshid Mas'ud al-Kashi\\
      % Arabic: غياث الدين جمشيد مسعود الكاشي
      }\\[3pt]
      \textbf{Dates:} c.\ 1380 -- 1429 CE\\
      \textbf{Origin:} Kashan (modern Iran);\\
      worked at the Samarkand observatory under Ulugh Beg
    \end{personbox}
    \vspace{0.3em}
    \begin{itemize}
      \item Calculated $\pi$ to \textbf{16 decimal places} --- a world record that stood for $\sim$180 years
      \item Invented an iterative algorithm for computing $n$th roots
      \item Wrote \textit{The Key to Arithmetic} (1427), a comprehensive mathematical encyclopedia
      \item Systematic use of \textbf{decimal fractions}
    \end{itemize}

    \column{0.42\textwidth}
    % Pi digits display
    \visualplaceholder[5cm]{Diagram}

    \vspace{0.6em}
    \visualplaceholder[3.5cm]{The Ulugh Beg observatory\\in Samarkand (modern photo)}
  \end{columns}

  \note{
    \textbf{Talking point:} ``He computed pi to 16 decimal places --- by hand. That record
    was not broken until the 1590s.''

    Al-Kashi worked at the observatory of Ulugh Beg, the Timurid ruler who was himself
    an accomplished astronomer. The observatory in Samarkand was one of the finest in
    the world.

    \verifytag{Exact number of decimal places sometimes reported as 16 or 17. Verify
    the standard claim before presenting.}
  }
\end{frame}

% ---- Slide 50: Islamic Mathematics Legacy ----
\begin{frame}{Islamic Mathematics: Legacy and Transmission}
  \begin{columns}[T]
    \column{0.45\textwidth}
    {\small Islamic mathematicians did not merely ``preserve'' Greek knowledge --- they \textbf{transformed} and \textbf{extended} it:}
    \vspace{0.3em}
    \begin{itemize}\setlength\itemsep{2pt}
      \item[\textcolor{islamicGreen}{\faCheck}] {\small\textbf{Invented} algebra as a discipline}
      \item[\textcolor{islamicGreen}{\faCheck}] {\small\textbf{Advanced} trigonometry to independence}
      \item[\textcolor{islamicGreen}{\faCheck}] {\small\textbf{Developed} computational methods}
      \item[\textcolor{islamicGreen}{\faCheck}] {\small\textbf{Pushed} toward symbolic thinking}
      \item[\textcolor{islamicGreen}{\faCheck}] {\small\textbf{Extended} number theory, optics, geodesy}
    \end{itemize}

    \vspace{0.4em}
    \begin{insightbox}
      {\footnotesize The Islamic mathematical tradition was essential to the development of Renaissance mathematics, Enlightenment science, and the modern world.}
    \end{insightbox}

    \column{0.52\textwidth}
    % Knowledge transmission map
    \visualplaceholder[5cm]{Simplified map outline (just positioning)}
  \end{columns}

  \vspace{0.3em}
  \centering
  {\small\textit{``Let us now see what happened when these ideas reached medieval Europe\ldots''}}

  \note{
    \textbf{Key message:} Emphasize ORIGINAL contributions, not just preservation.

    The transmission happened through specific contact points:
    \begin{itemize}
      \item Toledo (Reconquista --- scholars translated Arabic texts into Latin)
      \item Sicily (Norman conquest --- multilingual court)
      \item The Crusader states (limited but real intellectual exchange)
    \end{itemize}

    End of Session 2. Natural break point.
  }
\end{frame}
